\chapter*{Introduction}
\addcontentsline{toc}{chapter}{Introduction}
The neural network is a supervised machine learning model inspired by the biological neural tissue. Due to their ability to capture complex patterns, they are widely used in areas such as image recognition, language processing, and bioinformatics. However, their performance can deteriorate when datasets contain irrelevant or redundant features. High-dimensional datasets require longer training times and are harder to interpret, which is a problem often referred to as the \textit{curse of dimensionality}. To address this problem, feature selection methods are commonly applied.

Feature selection is the process of identifying the most relevant features in a dataset while removing those that do not contribute significantly to predictive performance. By reducing the number of input variables, we can improve model accuracy, speed up training, and make the results easier to interpret. 

The main goal of this thesis is to explore different feature selection methods in the context of neural networks. The work includes a theoretical comparison of these approaches and a practical demonstration on a selected dataset, showing how different methods influence model performance.

The structure of the thesis is as follows: Section 1 provides a theoretical overview of neural networks and information theory. Section 2 focuses on feature selection methods and their taxonomy. Section 3 presents the research methodology and experimental results. The Conclusion summarizes the main findings and their implications.

Overall, this work aims to highlight how appropriate feature selection can reduce computational complexity and improve the efficiency of neural network models, particularly in high-dimensional settings.